%%
%% This is file `pagecolor-example.tex',
%% generated with the docstrip utility.
%%
%% The original source files were:
%%
%% pagecolor.dtx  (with options: `example')
%% 
%% This is a generated file.
%% 
%% Project: pagecolor
%% Version: 2025-01-28 v1.2d
%% Info: Provides thepagecolor
%% 
%% Copyright (C) 2011 - 2025 by
%%     H.-Martin Muench <Martin dot Muench at Uni-Bonn dot de>
%% 
%% This work may be distributed and/or modified under the
%% conditions of the LaTeX Project Public License, either
%% version 1.3c of this license or (at your option) any later
%% version. This version of this license is in
%%    https://www.latex-project.org/lppl/lppl-1-3c.txt
%% and the latest version of this license is in
%%    https://www.latex-project.org/lppl.txt
%% and version 1.3c or later is part of all distributions of
%% LaTeX version 2005-12-01 or later.
%% 
%% This work has the LPPL maintenance status "maintained".
%% 
%% The Current Maintainer of this work is H.-Martin Muench.
%% 
%% This work consists of the main source file pagecolor.dtx,
%% the README, and the derived files
%%    pagecolor.sty, pagecolor.pdf,
%%    pagecolor.ins, pagecolor.drv,
%%    pagecolor-example.tex, pagecolor-example.pdf.
%% 
%% In memoriam
%%  Claudia Simone Barth + 1996-01-30
%%  Tommy Muench         + 2014-01-02
%%  Hans-Klaus Muench    + 2014-08-24
%% 
\NeedsTeXFormat{LaTeX2e}[2024-11-01]
\documentclass[british]{article}[2024/06/29]% v1.4n Standard LaTeX document class
\usepackage[extension=pdf,%
 plainpages=false,%
 pdfpagelabels=true,%
 hyperindex=false,%
 pdflang={en},%
 pdftitle={pagecolor package example},%
 pdfauthor={H.-Martin Muench},%
 pdfsubject={Example for the pagecolor package},%
 pdfkeywords={LaTeX, pagecolor, thepagecolor, page color, page colour},%
 pdfview=Fit,pdfstartview=Fit,%
 pdfpagelayout=SinglePage%
]{hyperref}[2024-11-05]% v7.01l Hypertext links for LaTeX

\usepackage[x11names]{xcolor}[2024/09/29]% v3.02 LaTeX color extensions (UK)
 % The xcolor package would not be needed for just using the base colors.
 % The color package would be sufficient for that.

 % \usepackage[cam,center,a3]{crop}[2017/11/19]% v1.10

\usepackage[pagecolor={LightGoldenrod1},%
  nopagecolor={none}]{pagecolor}[2025-01-28]% v1.2d Provides thepagecolor (HMM)

\usepackage{afterpage}[2023/07/04]% v1.08 After-Page Package (DPC)
 % The afterpage package is generally not needed,
 % but the |\newpagecolor{somecolor}\afterpage{\restorepagecolor}|
 % construct shall be demonstrated.

\usepackage{lipsum}[2021-09-20]% v2.7 150 paragraphs of Lorem Ipsum dummy text
 % The lipsum package is generally not needed,
 % but some blind text is needed for the example.

\listfiles
\begin{document}
\pagenumbering{arabic}
\section*{Example for pagecolor}

This example demonstrates the use of package\newline
\textsf{pagecolor}, v1.2d as of 2025-01-28 (HMM).\newline
The used options were\newline
\verb|pagecolor={LightGoldenrod1}|\newline
(\verb|pagecolor={none}| would be the default), and\newline
\verb|nopagecolor={none}| (which is the default).

\noindent For more details please see the documentation!

\noindent The current page (background) color is\newline
\verb|\thepagecolor|\ =\ \thepagecolor \newline
(and \verb|\thepagecolornone|\ =\ \thepagecolornone ,
which would only be different from \verb|\thepagecolor|,
when the page color would be \verb|none|).

\newpage
\pagecolor{rgb:-green!40!yellow,3;green!40!yellow,2;red,1}

{\color{white} The current page (background) color is\newline
\verb|\thepagecolor|\ =\ \thepagecolor .}

{\color{\thepagecolor} And that makes this text practically invisible.}

{\color{white} Which made the preceding line of text practically
invisible, but it can be copied and pasted.}

\newpage
\newpagecolor{red}

This page uses \verb|\newpagecolor{red}|.

\newpage
\restorepagecolor

{\color{white}And this page uses \verb|\restorepagecolor| to restore
the page color to the value it had before the red page.}

\newpage
\pagecolor{none}

This page uses \verb|\pagecolor{none}|. If the \verb|\nopagecolor|
command is known, the page color is now
\verb|none| (because option \verb|nopagecolor={none}|), otherwise
\verb|white| (or the color given with option \verb|nopagecolor={...}|):
\newline
\verb|\thepagecolor|\ =\ \thepagecolor\ and
\verb|\thepagecolornone|\ =\ \thepagecolornone .

\newpage
\restorepagecolor

{\color{white}\verb|\restorepagecolor| restored the page color again.}

\newpage
\pagecolor{green}

This page is green due to \verb|\pagecolor{green}|.

\newpage
\newpagecolor{blue}\afterpage{\restorepagecolor}

{\color{white}\verb|\newpagecolor{blue}\afterpage{\restorepagecolor}|%
\newline
was used here, i.\,e.~this page is blue, and the next one will
automatically have the same page color before it was changed to blue
here (i.\,e. green).}

\smallskip
{\color{red}\textbf{\lipsum[1-11]}}
\bigskip

The page color was changed back at the end of the page --
in mid-sentence!

\newpage
\backgroundpagecolor{pink}

When activating the loading of the crop package in the preamble of this
document, \verb|\backgroundpagecolor{<|\textit{some color}\verb|>}|
changes the color of the background/outer/physical page.

\newpage
\newbackgroundpagecolor{blue}

Analogous to \verb|\newpagecolor{...}| and \verb|\restorepagecolor|,
for the background/outer/physical page
\verb|\newbackgroundpagecolor{<|\textit{some color}\verb|>}| and\linebreak
\verb|\restorebackgroundpagecolor| are provided.

Here \verb|\newbackgroundpagecolor{blue}| colored that
background/outer/physical page in blue (if crop is used).

\newpage
\restorebackgroundpagecolor

And here the pink color of the background/outer/physical page
was restored by \verb|\restorebackgroundpagecolor| (if crop is used).

\end{document}
\endinput
%%
%% End of file `pagecolor-example.tex'.
